\documentclass[10pt,a4paper]{amsart}
\usepackage[margin=19mm]{geometry}
\usepackage{amsmath,amssymb,mathtools}
\usepackage[scr=rsfs,cal=euler]{mathalfa}
\usepackage{xfrac}
\usepackage[unicode,hidelinks,bookmarks=false]{hyperref}
\newcommand{\ie}{i.e.\ }
\newcommand{\cf}{cf.\ }
\newcommand{\resp}{respectively\ }
\newcommand{\Subsection}{\subsection}
\providecommand{\nobreakdash}{\nobreak\hbox{-}}
\setlength{\emergencystretch}{2em}
\multlinegap0pt
\usepackage{bm}
\usepackage{amssymb}
\usepackage{xcolor}
\usepackage[shortlabels]{enumitem}
\newtheorem{lemma}{Lemma}
\title{Asymptotic Properties of \\ Generalized Elephant Random Walks}
\author{Krishanu Maulik, Parthanil Roy, Tamojit Sadhukhan}
\date{Source: arXiv:2406.19383v2 (CC BY 4.0)}
\begin{document}
\maketitle

\noindent\emph{Source and licence.} Asymptotic Properties of \\ Generalized Elephant Random Walks. Version arXiv:2406.19383v2 (CC BY 4.0), distributed by arXiv under the Creative Commons Attribution 4.0 International licence, http://creativecommons.org/licenses/by/4.0/, as recorded in the arXiv licence metadata for this exact version and shown on the abstract page https://arxiv.org/abs/2406.19383v2. Full licence notice: \texttt{licenses/arxiv-2406.19383-CC-BY-4.0.txt}.\par\smallskip

\noindent\textit{Editorial scope.} Counted unit: the article's lemma lem-noiseanddrift2 (source0.tex lines 1198--1235). The article's complete original proof of it is delivered with it (source0.tex lines 1238--1301, one \texttt{proof} environment of 64 lines). No mathematics is written, repaired or re-derived: the statement and the proof are the article's own lines, in the article's own notation, and no retained line is substituted or re-flowed.  Retained uncounted as the source-local context the proof needs: lines 519--524; lines 1186--1188.  Nothing was omitted from any retained span, so the package carries no omission record.  No retained line contains a citation, so the wrapper carries no bibliography and no bibliography entry.  No retained line contains a figure environment, an image-inclusion command or drawing code, so no image is delivered and the package declares no asset dependency.  Editorial formatting only, in addition: an AMS \texttt{amsart} preamble that restates the article's own declarations and macros used by the retained lines, namely the environments \texttt{lemma} on separate counters, together with the standard packages those lines need.  The retained environments print in this document as Lemma 1 in order of appearance, whereas the article prints the same items with the numbering of its own full document.  The complete native source of the article as distributed in the arXiv e-print for this exact version is bundled unchanged as \texttt{sources/161476/source0.tex}.
\begin{align}\label{eq12}
           \boldsymbol{\widetilde{X}}_{n+1} = \begin{cases}
               \boldsymbol{Y}^{({\pi}_1)}_{n+1} &\text{ with probability }   \mathcal{P}_1\left(\frac{\boldsymbol{\widetilde{S}}_n}{n}\right), \\ \boldsymbol{Y}^{({\pi}_2)}_{n+1} &\text{ with probability }   \mathcal{P}_2\left(\frac{\boldsymbol{\widetilde{S}}_n}{n}\right), \\ \vdots \\ \boldsymbol{Y}^{({\pi}_{r-1})}_{n+1} &\text{ with probability }   \mathcal{P}_{r-1}\left(\frac{\boldsymbol{\widetilde{S}}_n}{n}\right), \\ 
               \boldsymbol{Y}^{({\pi}_r)}_{n+1} &\text{ with probability } 1-\sum_{j=1}^{r-1}\mathcal{P}_j\left(\frac{\boldsymbol{\widetilde{S}}_n}{n}\right).
           \end{cases}
\end{align}

\begin{align}\label{eq-sa2}
   \boldsymbol{\Gamma}_{n+1} = \boldsymbol{\Gamma}_{n} - \alpha_n\left(\boldsymbol{\gamma}(\boldsymbol{\Gamma}_{n}) + \boldsymbol{e}_{n+1}\right).
\end{align}

\begin{lemma}[]\label{lem-noiseanddrift2}
{Assume that $\mathbb{E}\left(\|\boldsymbol{Y}_1\|^2\right) < \infty$. Then the random noise $\left\{\boldsymbol{e}_{n}\right\}_{n\geq 2}$ of the %
process $\left(\boldsymbol{\Gamma}_{n}, \mathcal{G}_{n}\right)$, given by~\eqref{eq-sa2}, satisfies the following properties. 
\begin{enumerate}
\item{}\label{item2} The conditional distribution of $\boldsymbol{\widetilde{X}}_{n+1}$ given $\mathcal{G}_{n}$ is a measurable function of $\boldsymbol{\Gamma}_{n}$ alone and so the same holds for $\boldsymbol{e}_{n+1}$ as well. More precisely, for any bounded measurable function $g : D_s \to \mathbb{R}$, 
\begin{align}\label{eq:g}
   \mathbb{E}\left(g\left(\boldsymbol{\widetilde{X}}_{n+1}\right) \mid \mathcal{G}_{n}\right) = \sum_{i=1}^{r}\mathcal{P}_i\left(\boldsymbol{\Gamma}_{n}\right)\mathbb{E}\left(g\left(\boldsymbol{Y}^{(\pi_i)}\right)\right).
\end{align}
Moreover, $\left(\boldsymbol{e}_{n}\right)_{n \geq 2}$ is an $L^2$ martingale-difference sequence with respect to the filtration $\left(\mathcal{G}_{n}\right)_{n \geq 2}$. 

Define the functions $\sigma^2$ and $\boldsymbol{\Sigma}$ on ${D}_s$ as: \begin{align}\label{eq:lem-Sigma}
\sigma^2\left(\boldsymbol{x}\right) := \mathbb{E}\left(\|\boldsymbol{e}_{n+1}\|^2 \mid \boldsymbol{\Gamma}_{n} = \boldsymbol{x}\right), \quad \boldsymbol{\Sigma}\left(\boldsymbol{x}\right) := \mathbb{E}\left(\boldsymbol{e}_{n+1}\boldsymbol{e}^{\top}_{n+1} \mid \boldsymbol{\Gamma}_{n} = \boldsymbol{x}\right).
\end{align}
Then we have
\begin{align*}%
\boldsymbol{\Sigma}\left(\boldsymbol{x}\right) &= \sum_{i=1}^{r}\mathcal{P}_i\left(\boldsymbol{x}\right)\boldsymbol{\Sigma}^{(\pi_i)}
  -  \boldsymbol{H}\left(\boldsymbol{x}\right)\left(\boldsymbol{H}\left(\boldsymbol{x}\right)\right)^{\top},
\intertext{and}
\sigma^2\left(\boldsymbol{x}\right) &= \sum_{i=1}^{r}\mathcal{P}_i\left(\boldsymbol{x}\right)\operatorname{tr}\boldsymbol{\Sigma}^{(\pi_i)} - \left(\boldsymbol{H}\left(\boldsymbol{x}\right)\right)^{\top}\boldsymbol{H}\left(\boldsymbol{x}\right).
\end{align*}
For some constant $C > 0$ and for all $\boldsymbol{x} \in {D}_s$, we also have \begin{align*}%
\mathbb{E}\left(\|\boldsymbol{\gamma}(\boldsymbol{\Gamma}_n) + \boldsymbol{e}_{n+1}\|^2 \mid \boldsymbol{\Gamma}_n = \boldsymbol{x}\right) \leq \|\boldsymbol{\gamma}(\boldsymbol{x})\|^2 + \sigma^2(\boldsymbol{x}) \leq C(1 + \|\boldsymbol{x}\|^2).
\end{align*}

Additionally, if $\boldsymbol{\Gamma}_n \stackrel{a.s.}{\to} \boldsymbol{x}_0$ for some $\boldsymbol{x}_0 \in {D}_s$ and $\boldsymbol{\mathcal{P}}$ is continuous at $\boldsymbol{x}_0$, then
\begin{align}\label{eq:lem-Sigmaconv}
\sigma^2\left(\boldsymbol{\Gamma}_{n}\right) \stackrel{a.s.}{\to} \sigma^2\left(\boldsymbol{x}_0\right), \quad \boldsymbol{\Sigma}\left(\boldsymbol{\Gamma}_{n}\right) \stackrel{a.s.}{\to} \boldsymbol{\Sigma}\left(\boldsymbol{x}_0\right).
\end{align}
\item{}\label{item3} The following Lindeberg condition is satisfied: for all $\delta > 0$,
    \begin{align}\label{eq:lind}
    \frac{1}{n}\sum_{k=1}^{n}\mathbb{E}\left( \|\boldsymbol{e}_{k+1}\|^2 \mathbf{1}\left\{\|\boldsymbol{e}_{k+1}\| \geq \delta\sqrt{n}\right\}\mid \mathcal{G}_{k}\right) \stackrel{a.s.}{\to} 0.
    \end{align}
If we additionally have $\mathbb{E}\left(\|\boldsymbol{Y}_1\|^{2+\epsilon}\right) < \infty$ for some $\epsilon > 0$, then
\begin{align}\label{eq:sup}
        \sup_{n \geq 1} \mathbb{E}\left( \|\boldsymbol{e}_{n+1}\|^{2+\epsilon} \mid \mathcal{G}_{n}\right) < \infty.
\end{align}
\end{enumerate}}
\end{lemma}

\begin{proof} 
{%
The assertion that the conditional distribution of $\boldsymbol{\widetilde{X}}_{n+1}$ given $\mathcal{G}_{n}$ is a measurable function of $\boldsymbol{\Gamma}_{n}$ alone and \eqref{eq:g} follow from \eqref{eq12}. Next observe that $\|\boldsymbol{\widetilde{X}}_{n+1}\| \leq \|\boldsymbol{Y}_{n+1}\|$ and for any $\boldsymbol{x} \in {D}_s$,
\begin{align}\label{eq:h-ineq}
    \|\boldsymbol{H}(\boldsymbol{x})\| &= \left\|\sum_{i=1}^{r}\mathcal{P}_i(\boldsymbol{x})\boldsymbol{\mu}^{(\pi_i)}\right\| \leq \|\boldsymbol{\mu}\| \sum_{i=1}^{r} \mathcal{P}_i(\boldsymbol{x}) = \|\boldsymbol{\mu}\|. 
\end{align}
Thus, for any $n \geq 1$,
\begin{align}\label{eq:e-bound}
    \|\boldsymbol{e}_{n+1}\| &\leq \left\|\boldsymbol{H}\left({\boldsymbol{\Gamma}_n}\right)\right\| + \|\boldsymbol{\widetilde{X}}_{n+1}\| \leq \|\boldsymbol{\mu}\| + \|\boldsymbol{Y}_{n+1}\|,
\end{align}
and so
\[\mathbb{E}\left(\|\boldsymbol{e}_{n+1}\|\right) \leq \|\boldsymbol{\mu}\| +  \mathbb{E}\left(\|\boldsymbol{Y}_1\|\right)< \infty.\]}

Also, from 
\eqref{eq12}, it follows that 
\begin{align*}
    \mathbb{E}\left(\boldsymbol{e}_{n+1} \mid \mathcal{G}_{n}\right) = \mathbb{E}\left(\boldsymbol{H}({\boldsymbol{\Gamma}_n}) - \boldsymbol{\widetilde{X}}_{n+1} \Big| \boldsymbol{\widetilde{S}}_{n}\right) = \boldsymbol{H}({\boldsymbol{\Gamma}_n}) -  \sum_{i=1}^{r}\mathcal{P}_i(\boldsymbol{x})\boldsymbol{\mu}^{(\pi_i)} 
    = 0.
\end{align*}
Thus 
$\left(\boldsymbol{e}_{n}\right)_{n \geq 2}$ is a martingale-difference sequence with respect to $\left(\mathcal{G}_{n}\right)_{n \geq 2}$.
Using~\eqref{eq:h-ineq}, 
for any $n \geq 1$, we get
\begin{align}\label{eq:e-bound2}
    \|\boldsymbol{e}_{n+1}\|^2 &\leq 2\left(\left\|\boldsymbol{H}\left({\boldsymbol{\Gamma}_n}\right)\right\|^2 + \|\boldsymbol{\widetilde{X}}_{n+1}\|^2\right)
   \leq 2\left(\|\boldsymbol{\mu}\|^2 + \|\boldsymbol{Y}_{n+1}\|^2\right),
\end{align}
which makes the martingale difference $L^2$-bounded.

{From \eqref{eq12}, we further have
\begin{align*}
\boldsymbol{\Sigma}\left(\boldsymbol{x}\right) = \mathbb{E}\left(\boldsymbol{e}_{n+1}\boldsymbol{e}^{\top}_{n+1} \mid \boldsymbol{\Gamma}_{n} = \boldsymbol{x}\right) &= \mathbb{E}\left( \boldsymbol{\widetilde{X}}_{n+1}\boldsymbol{\widetilde{X}}^{\top}_{n+1} \Big| \boldsymbol{\Gamma}_{n} = \boldsymbol{x}\right) - 
\boldsymbol{H}\left(\boldsymbol{x}\right)\left(\boldsymbol{H}\left(\boldsymbol{x}\right)\right)^{\top}
        \\ &= \sum_{i=1}^{r}\mathcal{P}_i\left(\boldsymbol{x}\right)\boldsymbol{\Sigma}^{(\pi_i)}
        -  \boldsymbol{H}\left(\boldsymbol{x}\right)\left(\boldsymbol{H}\left(\boldsymbol{x}\right)\right)^{\top}.
\end{align*}
As $\mathbb{E}\left(\|\boldsymbol{e}_{n+1}\|^2 \mid \boldsymbol{\Gamma}_{n} = \boldsymbol{x}\right) = \operatorname{tr} \mathbb{E}\left(\boldsymbol{e}_{n+1}\boldsymbol{e}^{\top}_{n+1} \mid \boldsymbol{\Gamma}_{n} = \boldsymbol{x}\right)$,
we have 
\begin{align*}%
\sigma^2\left(\boldsymbol{x}\right) &= \operatorname{tr} \boldsymbol{\Sigma} (\boldsymbol{x}) 
 = \sum_{i=1}^{r}\mathcal{P}_i\left(\boldsymbol{x}\right)\operatorname{tr}\boldsymbol{\Sigma}^{(\pi_i)} - \left(\boldsymbol{H}\left(\boldsymbol{x}\right)\right)^{\top}\boldsymbol{H}\left(\boldsymbol{x}\right).
  \end{align*}}
Further, 
$\sigma^2\left(\boldsymbol{x}\right) \leq \sum_{i=1}^{r}\mathcal{P}_i\left(\boldsymbol{x}\right)\operatorname{tr}\boldsymbol{\Sigma}^{(\pi_i)} 
  \leq \mathbb{E}\left(\|\boldsymbol{Y}_1\|^2\right)$. %

{Also, for all $\boldsymbol{x} \in {D}_s$,
\begin{align*}
    \|\boldsymbol{\gamma}\left(\boldsymbol{x}\right)\|^2 = \|\boldsymbol{x}-\boldsymbol{H}(\boldsymbol{x})\|^2 \leq 2\left(\|\boldsymbol{x}\|^2 + \|\boldsymbol{H}(\boldsymbol{x})\|^2\right) \leq 2\|\boldsymbol{x}\|^2 + 2 \|\boldsymbol{\mu}\|^2.
\end{align*}
Thus, for all $\boldsymbol{x} \in {D}_s$, we get a constant $C > 0$ such that
\[
  \mathbb{E}\left(\|\boldsymbol{\gamma}(\boldsymbol{\Gamma}_n) + \boldsymbol{e}_{n+1}\|^2 \mid \boldsymbol{\Gamma}_n = \boldsymbol{x}\right) \leq \|\boldsymbol{\gamma}\left(\boldsymbol{x}\right)\|^2 +  \sigma^2\left(\boldsymbol{x}\right) \leq C(1+\|\boldsymbol{x}\|^2).
\]}%
{The fact that $\boldsymbol{\Gamma}_n \stackrel{a.s.}{\to} \boldsymbol{x}_0$ and the continuity of $\boldsymbol{\mathcal{P}}$ at $\boldsymbol{x}_0$ imply $\boldsymbol{\Sigma}\left(\boldsymbol{\Gamma}_n\right) \stackrel{a.s.}{\to} \boldsymbol{\Sigma}\left(\boldsymbol{x}_0\right)$. Since $\operatorname{tr}$ is continuous operator, $\boldsymbol{\Sigma}\left(\boldsymbol{\Gamma}_n\right) \stackrel{a.s.}{\to} \boldsymbol{\Sigma}\left(\boldsymbol{x}_0\right)$ implies ${\sigma^2}\left(\boldsymbol{\Gamma}_n\right) \stackrel{a.s.}{\to} {\sigma^2}\left(\boldsymbol{x}_0\right)$.}

Now fix $\delta > 0$. Then, using \eqref{eq:e-bound},~\eqref{eq:e-bound2} and finite second moment of $\|\boldsymbol{Y}_1\|$, we get,
\begin{align*}%
&\frac{1}{n}\sum_{k=1}^{n}\mathbb{E}\left( \|\boldsymbol{e}_{k+1}\|^2 \mathbf{1}\left\{\|\boldsymbol{e}_{k+1}\| \geq \delta\sqrt{n}\right\}\mid \mathcal{G}_{k}\right) \nonumber
\\ \leq& \frac{2}{n}\sum_{k=1}^{n}\mathbb{E}\left( \left(\|\boldsymbol{\mu}\|^2 + \|\boldsymbol{Y}_{k+1}\|^2\right) \mathbf{1}\left\{\|\boldsymbol{\mu}\| + \|\boldsymbol{Y}_{k+1}\| \geq \delta\sqrt{n}\right\}\mid \mathcal{G}_{k}\right) \nonumber
\\ =& 2\mathbb{E}\left( \left(\|\boldsymbol{\mu}\|^2 + \|\boldsymbol{Y}_{1}\|^2\right) \mathbf{1}\left\{\|\boldsymbol{\mu}\| + \|\boldsymbol{Y}_{1}\| \geq \delta\sqrt{n}\right\}\right) \to 0. \nonumber
\end{align*}
{The assumption that $\mathbb{E}\left(\|\boldsymbol{Y}_1\|^{2+\epsilon}\right) < \infty$, combined with \eqref{eq:e-bound}, prove \eqref{eq:sup}.}
\end{proof}

\end{document}
